\paragraph{Two further positive cases, including a non-$F$-pure test.}

\begin{proposition}\label{prop:DE-E2-flat-charzero}
Let $k$ be a field of characteristic zero and let $X$ be locally of finite type over $k$.  Restriction gives an equivalence between quasi-coherent crystals whose values are flat modules and ordinary stratified quasi-coherent $\OO_X$-modules that are flat over $\OO_X$.  No finite-generation hypothesis is imposed.
\end{proposition}
\begin{proof}
A flat stratified module defines a derived object on the intrinsic completed nerve: every ordinary pullback appearing in its stratification agrees with its derived pullback and is flat over the corresponding finite diagonal neighbourhood.  The relevant mapping spaces are $0$-truncated, so the ordinary cocycle identities supply the derived coherences, exactly as for vector bundles.  Gaitsgory--Rozenblyum \cite[Proposition~3.4.3]{GaitsgoryRozenblyum2014} therefore gives a derived crystal $\mathcal F$.  We check that all its values are ordinary flat modules.

The needed nilpotent descent statement holds without any Noetherian or boundedness assumption.  Let $B$ be a ring, let $I^N=0$, and suppose $K\in D(B)$ has derived reduction $K\otimes_B^{\mathbf L}B/I=M[0]$, with $M$ flat over $B/I$.  For every $B/I$-module $Q$, associativity gives
\[
 K\otimes_B^{\mathbf L}Q\simeq M\otimes_{B/I}Q[0].
\]
Filter an arbitrary $B$-module $P$ by $I^jP$.  Its finitely many successive quotients are $B/I$-modules, so $K\otimes_B^{\mathbf L}P$ is concentrated in degree zero.  Taking $P=B$ first shows that $K$ is an ordinary module; varying $P$ then proves its flatness.  Apply this to each infinitesimal test of $X$.  Thus $\mathcal F$ is the required flat ordinary crystal.

Conversely, a flat ordinary crystal has derived Cartesian transition maps, since its values are flat.  The smooth-ambient extension to arbitrary classical affine tests used in Proposition~\ref{prop:DE-E2-coherent} applies verbatim, with flat modules in place of vector bundles.  The functorial construction also applies to morphisms.  Corollary~2.2.2 and Proposition~3.4.3 of the same source give full faithfulness and show that the two constructions are inverse.
\end{proof}

\begin{lemma}\label{lem:DE-Fdiv-universal-homeomorphism}
Over a perfect field of characteristic $p$, pullback along a finite universal homeomorphism between reduced schemes locally of finite type induces an equivalence on the categories of $F$-divided quasi-coherent modules, with arbitrary quasi-coherent coefficients.
\end{lemma}
\begin{proof}
The assertion is local on the target.  For an affine such morphism write $A\subset B$; injectivity follows because the schemes are reduced and the map is surjective.  By \cite[Lemma~29.47.9, Tag~0CNF]{StacksProject}, every $b\in B$ has a $p$-power in $A$.  Finitely many algebra generators of $B$ therefore give one $q=p^N$ with $B^q\subset A$.  The homomorphism $g^\#:B\to A$, $b\mapsto b^q$, then satisfies
\[
             f\circ g=F_X^N,\qquad g\circ f=F_Y^N.
\]
Here we use absolute Frobenius, so $g$ may be semilinear over the ground field.  Equivalently one may insert the corresponding relative Frobenius twists.  These descriptions agree because the ground field is perfect.

On $F$-divided systems, pullback by absolute Frobenius is an equivalence, with inverse the shift $(M_0,M_1,\ldots)\mapsto(M_1,M_2,\ldots)$; the structural isomorphisms give both inverse identities.  Consequently both composites $g^*f^*$ and $f^*g^*$ are equivalences, and $f^*$ is an equivalence.  Its fully faithful local inverses glue by quasi-coherent descent, proving the assertion.
\end{proof}

\begin{proposition}\label{prop:DE-E2-cusps}
Let $k$ be a perfect field of characteristic $p$, and let $X$ be an integral Cohen--Macaulay $k$-variety admitting a finite universal homeomorphism $f:Y\to X$ with $Y$ smooth.  Restriction is an equivalence between quasi-coherent crystals on $(X/k)_{\rm inf}$ and ordinary stratified quasi-coherent $\OO_X$-modules, with no finiteness or flatness requirement on the coefficients.

In particular this applies to the cusp $X=\Spec k[t^2,t^3]$, although this cusp is not $F$-pure for any prime $p$.
\end{proposition}
\begin{proof}
Ben-Zvi--Nevins \cite[Proposition~2.11, Theorem~3.11, Theorem~4.3 and Remark~4.5]{BenZviNevins2004} show that ordinary pullback $f^*$ is an equivalence on arbitrary quasi-coherent stratified modules in these hypotheses.  Their Remark~2.7 takes tensor products in the pro-category, so the stratifications used there are precisely the compatible ordinary stratifications on all finite diagonal neighbourhoods used here; no interchange with a completed tensor product is involved.  The good Cohen--Macaulay hypotheses needed for their Theorems~3.11 and~4.3 follow from smoothness of $Y$ and Proposition~2.11.

By Lemma~\ref{lem:DE-crystal-Fdiv} and Lemma~\ref{lem:DE-Fdiv-universal-homeomorphism}, $f^*$ is also an equivalence on quasi-coherent infinitesimal crystals.  The restriction functors commute with $f^*$.  The restriction functor for $Y$ is an equivalence by the smooth case, so the resulting commutative square of categories shows that the restriction functor for $X$ is an equivalence.

For the stated cusp, $k[t]$ is its finite normalization, and the map is bijective on geometric points, hence a universal homeomorphism.  The curve is a hypersurface and therefore Cohen--Macaulay.  To verify nonpurity directly, put $A=k[t^2,t^3]$ and $I=(t^2)$.  The class of $t^3$ in $A/I$ is nonzero, whereas its image under the Frobenius test map is zero:
\[
       t^3\notin(t^2),\qquad
       (t^3)^p=t^{3p}\in(t^{2p})=I^{[p]},
\]
because $t^p\in A$ for every $p\geq2$.  Thus the Frobenius map is not pure, and this positive case is not a consequence of Proposition~\ref{prop:DE-E2-Fpure}.
\end{proof}

These results leave the full relative E.2 assertion open: neither arbitrary nonflat coefficients in characteristic zero nor arbitrary relative flat morphisms are covered.  In characteristic $p$, the cusp calculation rules out that concrete non-$F$-pure family as a source of counterexamples.
